\section{Representation Learning and Autoencoders}

Dimensionality reduction aims to represent high-dimensional data using a smaller number of features while retaining relevant information. 
A widely used linear approach is principal component analysis (PCA), originally introduced by Pearson~\cite{Pearson1901}. 
PCA transforms the input into a set of orthogonal principal components, ordered by the amount of variance they explain. 
By retaining only a subset of these components, the dimensionality of the data can be reduced while preserving as much variance as possible.

Autoencoders (AEs) provide a learned approach to dimensionality reduction. 
An autoencoder consists of an encoder that maps the input to a lower-dimensional latent representation, and a decoder that reconstructs the original input from this representation. 
The model is commonly trained by minimizing the difference between the input and its reconstruction. 
Linear autoencoders are closely related to PCA and, when trained with a squared reconstruction objective under suitable conditions, can learn the same principal subspace, although their individual latent dimensions are not required to match the principal components~\cite{BaldiHornik1989}. 
Autoencoders can also use multiple nonlinear layers, allowing them to learn more complex representations than linear dimensionality reduction. 
Hinton and Salakhutdinov demonstrated the use of deep autoencoders for reducing the dimensionality of high-dimensional data and compared their learned representations with PCA~\cite{HintonSalakhutdinov2006}.

More generally, representation learning aims to learn useful features directly from the data instead of relying on manually defined representations~\cite{Bengio2013}. 
In autoencoders, the latent space provides such a representation by compressing information needed to reconstruct the input. 
The use of nonlinear transformations and multiple layers allows deep autoencoders to capture relationships that cannot be represented by a purely linear transformation. 
The structure of this latent space can therefore provide information about how the model represents patterns present in the original high-dimensional data.

Autoencoders can also be built using convolutional layers instead of fully connected layers. 
Convolutional autoencoders apply shared filters across the input, allowing local patterns to be learned while reducing the number of independent model parameters. 
Masci et al.\ introduced stacked convolutional autoencoders for learning hierarchical representations from data~\cite{Masci2011}. 
Convolutional architectures therefore provide an alternative to fully connected autoencoders when local structure in the input may contain useful information.

A further extension is the variational autoencoder (VAE), which introduces a probabilistic latent representation. 
Rather than encoding an input directly into a single latent vector, a VAE learns the parameters of a distribution from which the latent representation is sampled. 
Kingma and Welling introduced this approach together with the reparameterization method that allows the model to be trained through standard gradient-based optimisation~\cite{KingmaWelling2014}. 
The training objective combines reconstruction with a Kullback–Leibler (KL) divergence term that regularizes the latent distribution towards a chosen prior. 
This encourages a more structured and continuous latent space while still allowing the input to be reconstructed from its lower-dimensional representation.